\chapter{Contribution: Automated Vulnerability Remediation}
\label{ch:remediation}

% =====================================================================
\section{Introduction}
\label{rem:intro}
% =====================================================================

The preceding chapter established that locating injection-class
vulnerabilities in source code is, for the most part, a problem the
community already knows how to solve. Industrial static analysers trace
tainted data from an untrusted source to a dangerous sink and report the
offending flow with a precision that was unimaginable a decade ago. Yet a
report is not a repair. Once a finding is raised, somebody (almost always a
human) must still decide what the code was meant to do, separate the fixed
command skeleton from the variable data woven into it, and rewrite the call
so that the data can never again be mistaken for syntax. That last step is
where the security debt accumulates, and it is the subject of this chapter.

Injection vulnerabilities are a useful lens because they are at once
old, well understood, and stubbornly common. SQL injection (CWE-89),
OS-command injection (CWE-78), LDAP injection (CWE-90) and XPath injection
(CWE-643) all share a single root cause: untrusted data is concatenated
into a string that an interpreter later parses as a command
\cite{su2006essence,halfond2006classification}. The pattern has sat at or
near the top of the OWASP Top~Ten for the better part of twenty years
\cite{owasptop10}. It persists not because the defect is mysterious: any
competent engineer can explain it, but because the sheer volume of
vulnerable code dwarfs the human capacity to repair it correctly, one sink
at a time. We therefore frame remediation, rather than detection, as the
binding constraint on secure software delivery.

A natural reaction is to automate the fix the same way detection was
automated. The dominant contemporary recipe is to treat repair as a
translation task: collect pairs of vulnerable and fixed code, train a
sequence model, and let it propose an edit
\cite{fu2022vulrepair,chi2023seqtrans}. More recent variants prompt a large
language model, sometimes in a feedback loop with a test harness
\cite{xia2023chatrepair,jin2023inferfix,placeholder-llm-repair}. These
systems are impressive, and they undeniably fix some bugs. The difficulty
is more subtle than raw success rate. A learned model produces an edit that
\emph{looks} right; it offers no guarantee that the edit \emph{is} right.
In a security setting, that gap is not academic. A patch that compiles,
passes a smoke test, and silences the scanner while remaining
exploitable, and is arguably worse than no patch at all, because it retires the
ticket and manufactures false confidence. We return to this failure mode
repeatedly, since it motivates almost every design decision that follows.

This chapter develops \irsam{}, a remediation engine constructed on a
deliberately uncompromising premise. An automated repair tool, we argue,
must either emit a patch it can \emph{prove} safe against a precise,
machine-checkable criterion, or it must \emph{abstain}, declining to act and
explain why, in a form a downstream human or tool can consume. There is no
permitted third option in which the engine ships an edit it cannot justify.
This is a strong commitment, and it has a cost: \irsam{} repairs fewer
findings than a system willing to guess. The thesis of the chapter is that
the trade is worth making, and that the resulting guarantees are both
formally statable and empirically attainable.

\paragraph{Chapter objectives.}
\label{sec:objectives}
The work is organised around six concrete objectives, denoted R1 through R6
and revisited throughout the chapter:
\begin{enumerate}[label=\textbf{R\arabic*.},leftmargin=2.8em]
  \item \textbf{Soundness.} Every \emph{accepted} patch must satisfy a
  precise, decidable soundness criterion: attacker-controlled data can never
  reach the interpreter as syntax.
  \item \textbf{Sound abstention.} When no patch can be proven sound, the
  engine returns \bottom{} together with a machine-readable reason, rather
  than emitting an unverified edit.
  \item \textbf{Intent fidelity.} A patched program must behave identically
  to the original on benign, trusted inputs.
  \item \textbf{Detector-agnostic operation.} The engine consumes the output
  of any standard SAST tool through one unified finding schema.
  \item \textbf{Multi-interpreter, multi-language reach.} A single pipeline
  spans SQL, OS-command, LDAP and XPath sinks across Java and Python.
  \item \textbf{Auditability.} Every accepted patch carries a structural
  certificate that an independent reviewer or verifier can re-check.
\end{enumerate}

\paragraph{Chapter outline.}
The remainder of the chapter proceeds as follows.
Section~\ref{rem:problem} formalises the remediation problem, fixes the
threat model and adversarial assumptions, and states the hypotheses that the
rest of the work tests. Section~\ref{rem:design} presents the proposed
approach: the seven-stage pipeline, the Intent Abstract Model at its centre,
the binding function that turns intent into a host-language rewrite, and the
decidable soundness criterion that gates acceptance, closing with the
principal design decisions and their justification.
Section~\ref{rem:impl} descends into the implementation: the technology
stack, the unified finding schema, a stage-by-stage walkthrough of the
codebase, the declarative binder catalogs, the five validation gates, and
the Lean~4 mechanisation of the soundness argument.
Section~\ref{rem:setup} describes the experimental setup: environment,
corpora, metrics and baselines. Section~\ref{rem:results} reports and
analyses the results across a controlled tier and a realistic project tier,
including a comparison with representative baselines and an account of the
defects the evaluation surfaced. Section~\ref{rem:discussion} interprets the
findings, examines limitations and threats to validity, and sketches future
directions, before Section~\ref{rem:conclusion} concludes and ties the
contribution back to the objectives above.


% =====================================================================
\section{Problem Formulation}
\label{rem:problem}
% =====================================================================

Before a solution can be claimed sound, the problem it solves has to be
pinned down without ambiguity. This section does that work. We first set out
the threat model and the adversary we defend against, then give a compact
formal model of an injection sink, state the soundness criterion that
defines a correct repair, diagnose the three structural gaps that current
remediation leaves open, and finally crystallise the whole into a small set
of hypotheses. The formal vocabulary introduced here is reused, unchanged,
in every later section; it is, in effect, the contract between the theory
and the code.

\subsection{Threat model and adversarial assumptions}
\label{sec:threat}

We fix a host programming language $H$ (in this work Java or Python) and
an \emph{interpreter} $I=(L_I,\llbracket\cdot\rrbracket_I,P_I)$. Here $L_I
\subseteq \Sigma_I^{*}$ is the surface language of well-formed interpreter
commands (SQL queries, shell command lines, LDAP filters, XPath
expressions); $\llbracket\cdot\rrbracket_I$ is the operational semantics
that maps a command and a state to a successor state (database tuples, OS
process state, a directory service, an XML document); and $P_I$ is the
parser that turns a raw string into a concrete syntax tree or rejects it.
This three-part view of an interpreter, separating surface syntax,
semantics and parser, is what lets us reason about \emph{where} attacker
bytes are interpreted rather than merely \emph{whether} they are present.

The adversary is modelled through two disjoint value domains. The
\emph{trust domain} $D_T$ contains host values the developer controls:
string literals, configuration constants, and identifiers drawn from a
finite, known allow-list. The \emph{attack domain} $D_A$ contains values an
adversary can influence: HTTP parameters and headers, request bodies, file
contents, deserialised payloads, and any other channel crossing the trust
boundary. We assume $D_T \cap D_A = \emptyset$ modulo a conservative taint
analysis; that is, when in doubt a value is treated as attacker-controlled.
This assumption is deliberately pessimistic; it is safer to over-approximate
$D_A$ and abstain than to under-approximate it and ship a hole.

Three assumptions delimit what \irsam{} is responsible for. First, the
adversary controls data, not code: they can supply arbitrary inputs to
$D_A$ channels but cannot edit the program text, alter the binder catalog,
or tamper with the compiler. Second, the interpreter and its driver honour
the parameterised-query contract: a value delivered on the protocol channel
(for instance through \code{PreparedStatement.setString}) is bound as an
opaque datum and is never re-lexed as syntax. This is the standard
dialect-prepared-statement assumption, and it is the one external fact our
soundness argument leans on. Third, detection is a separate concern: we
assume some upstream tool localises sinks, and we make no claim about its
recall. The remediation guarantee must hold for whatever findings arrive,
and must never depend on the detector being complete.

\subsection{A formal model of an injection sink}
\label{sec:formal-model}

Let $e$ be a host expression that evaluates to a pair $(s,\sigma) \in
\Sigma_I^{*} \times \mathrm{State}$, and let $f$ be a sink API with a
tainted argument position $i$, so that the program point $f(\dots,e,\dots)$
ultimately invokes $\llbracket P_I(s)\rrbracket_I(\sigma)$. Following the
formal model that underpins the implementation, we call the program point
\emph{vulnerable} with respect to a malicious-behaviour set
$\mathrm{Mal}_I \subseteq \mathrm{State}'$ precisely when an attacker can
choose a $D_A$-controlled assignment to the free variables of $e$ such that
the resulting surface string parses ($P_I(s)\in L_I$) and its execution
lands in the malicious set:
\begin{equation}
  \exists\,\text{$D_A$-assignment}:\quad
  P_I(s)\in L_I \;\wedge\;
  \llbracket P_I(s)\rrbracket_I(\sigma)\in\mathrm{Mal}_I .
  \label{eq:vuln}
\end{equation}
In plain terms, the attacker bends the surface string into a shape the
interpreter accepts but the developer never intended. What counts as
malicious is fixed per interpreter: for SQL, any query shape outside the
developer's intended set; for the shell, any command outside it; and so on.

This phrasing already hints at the cure. The danger is not the presence of
untrusted data but its \emph{position}: the data sits inside the grammar,
on the side of the parser where bytes become tokens. Figure~\ref{fig:dataflow}
contrasts the two regimes. On the left, the value \code{name} is
concatenated into the query string and is therefore lexed as SQL; a crafted
input can close the quote and append syntax. On the right, the same value
is delivered through the driver's parameter channel: the parser sees only a
placeholder, and the value travels on a separate protocol path where it can
never be tokenised. The corrective transformation is thus not ``escape the
data more carefully'' but ``move the data out of the grammar.'' We will see
that expressing this move precisely is what forces a model of the
\emph{intended} command, and not merely of the byte stream.

\begin{figure}[ht]
\centering
\includegraphics[width=\linewidth]{figures/fig_dataflow.pdf}
\caption{The two regimes of an injection sink. \emph{Left:} in the
vulnerable program the untrusted value is concatenated into the query and
is lexed by the interpreter's parser, so a crafted input can inject syntax.
\emph{Right:} after remediation the value is bound on the driver's parameter
channel; the parser sees only a placeholder skeleton, and attacker bytes
travel on a protocol path that is never tokenised as code. The correct fix
relocates data from the grammar layer to the protocol layer.}
\label{fig:dataflow}
\end{figure}

\subsection{The soundness criterion}
\label{sec:soundness-crit}

Let $A$ denote the developer's \emph{intent}: the command they meant to
issue, parameterised by the positions where data flows in. Write
$\mathrm{Intent}(A)$ for the set of interpreter behaviours consistent with
that intent, $L_I$ for the well-formed commands, and $\mathrm{Mal}_I
\subseteq L_I$ for the malicious behaviours reachable by injection. A
remediation $p$ is \emph{sound} when, for every host state $\sigma$ and
every assignment of data positions to values drawn from $D_T \cup D_A$,
\begin{equation}
  \llbracket P_I(\mathrm{exec}(p,\sigma))\rrbracket_I(\sigma)
     \;\in\; \mathrm{Intent}(A) \;\subseteq\; L_I \setminus \mathrm{Mal}_I ,
  \label{eq:soundness}
\end{equation}
where $\mathrm{exec}(p,\sigma)$ is the command the patched program would
emit. The criterion says that the patched program's interpreter behaviour,
on \emph{any} input including adversarial ones, stays inside the intended
set and therefore outside the malicious set. The remediation problem, then,
is the following: given a finding that localises a sink, either produce a
patch $p$ for which Eq.~\eqref{eq:soundness} is \emph{provable}, or return
\bottom{}. Emitting a patch for which Eq.~\eqref{eq:soundness} cannot be
established is explicitly disallowed. This last clause is the formal
statement of the anti-false-fix commitment, and it is what distinguishes the
present formulation from one that merely maximises a fix rate.

It is worth dwelling on a subtlety. Eq.~\eqref{eq:soundness} quantifies over
\emph{all} inputs, including $D_A$ values at the value positions. Soundness
does not forbid the attacker from supplying a malicious string; it
guarantees that, however malicious the string, it is consumed as a single
value and not as syntax. Intent fidelity (objective R3) is the complementary
half: on $D_T$-only inputs the patched program reproduces the original
behaviour exactly. Together the two properties say that the patch changes
\emph{how} data reaches the interpreter without changing \emph{what} the
program computes on legitimate inputs.

\subsection{The remediation gap}
\label{sec:gap}

Why is this problem still open in practice? The discovery of injection flaws
is, as noted, largely a solved engineering task. Their remediation is not,
and three concrete deficiencies (which we label G1 through G3) explain the
shortfall. We state them as gaps in the current state of practice because
each one corresponds directly to an objective \irsam{} sets out to meet.

\paragraph{G1: No soundness contract.}
Learning-based and LLM-prompted repair systems emit edits that look correct
but carry no checkable invariant. ``Plausible'' is simply not the same as
``safe.'' A model may delete the offending line, paraphrase the query into
something subtly different, or insert a half-measure escape that a
determined attacker steps around. Crucially, there is no point in such a
pipeline at which the tool can be \emph{forced} to be right or, failing
that, to remain silent. The system is architecturally committed to
producing output, which is precisely the wrong commitment in a security
context.

\paragraph{G2: The detector--repairer impedance mismatch.}
Detectors emit \emph{locations}; safe repair needs \emph{intent}. A scanner
reports that tainted data reaches \code{executeQuery} on some line. It says
nothing about which substrings form the developer's fixed query skeleton and
which are the variable data that must be parameterised. A tool that operates
purely on the host abstract syntax tree therefore degrades to syntactic
pattern replacement and breaks the moment a command is assembled
dynamically, across statements, through helper methods, or with conditional
fragments. The information the repairer most needs is exactly the
information the detector does not provide.

\paragraph{G3: The false-fix failure mode.}
An incorrect patch that compiles and passes a cursory test is worse than no
patch. It closes the ticket, suppresses the alert, and ships a still
exploitable system wrapped in a false sense of safety. To our knowledge no
mainstream automated remediator treats ``I cannot prove this is safe'' as a
first-class, auditable outcome. The silent substitution of a guess for a
proof is the single most dangerous behaviour an automated security tool can
exhibit, and it is endemic to systems optimised for fix rate alone.

\subsection{Why escaping is the wrong fix}
\label{sec:escaping}

A recurring but structurally weak remedy deserves explicit rebuttal, because
it is so often proposed. The idea is to \emph{sanitise}: escape the
attacker's meta-characters inside the interpreter's lexer so that, for
instance, a quote in the input becomes an escaped quote in the query. This
keeps the untrusted data on the dangerous side of the parser and bets that
every meta-character has been escaped correctly, for every dialect, in every
context. The bet is fragile. Escaping rules differ across database engines,
across shell quoting modes, and across LDAP and XPath grammars; a single
missed case reopens the hole. Parameterised APIs take the opposite,
structurally sound route, already depicted in Figure~\ref{fig:dataflow}:
they relocate the value to the interpreter's protocol layer, where it is
bound once as an opaque datum and is never lexed as code
\cite{boyd2004sqlrand}. The distinction matters for our design because it
dictates \emph{what} the engine must reconstruct. To move data out of the
grammar, one must first know which parts of the command are grammar and
which are data; that is, one must recover the developer's intent.

\subsection{Hypotheses}
\label{sec:hypotheses}

The foregoing analysis can be distilled into three testable hypotheses that
the rest of the chapter examines. They are stated informally here and
revisited, against evidence, in Section~\ref{rem:results}.
\begin{description}[leftmargin=1.8em,style=nextline]
  \item[H1 (Provable repair is feasible).] For a useful fragment of
  real-world injection sinks, it is possible to reconstruct the developer's
  intent precisely enough to synthesise a patch whose soundness, in the
  sense of Eq.~\eqref{eq:soundness}, is decidable and re-checkable by an
  independent verifier.
  \item[H2 (Abstention is sound and frequent, not rare).] When intent cannot
  be reconstructed soundly, the engine can detect this and abstain with a
  precise reason; on realistic code such abstention is common, and treating
  it as a feature rather than a failure yields a system with no unsound
  patches.
  \item[H3 (Soundness need not collapse fix rate).] Insisting on provable
  repair does reduce coverage relative to a guess-freely baseline, but the
  reduction is bounded; on the supported fragment the functional fix rate
  remains substantially higher than that of baselines that do guess.
\end{description}

\subsection{Design premise: abstention as a feature}
\label{sec:premise}

Because a wrong fix is more dangerous than an unfixed-but-flagged finding
(G3), the rational design choice is a system that \emph{prefers silence to a
guess}. The single commitment, ``\bottom{} rather than an unsound
patch,'' reorients the entire architecture. Every stage of the pipeline is
permitted, and indeed required, to abstain the moment it cannot discharge
its local proof obligation, and each abstention carries a machine-readable
reason so that a human reviewer or a stronger downstream tool can take over
exactly where \irsam{} stops. This is the defining departure from generative
repair, where the model is structurally compelled to always output
something. We will see that making abstention first-class is not a
limitation grudgingly accepted but the very mechanism that makes the
soundness guarantee attainable.


% =====================================================================
\section{Proposed Approach and Architecture}
\label{rem:design}
% =====================================================================

Having fixed the problem, we can now describe the solution. This section
develops the architecture of \irsam{} from its guiding rationale down to its
component structure. The presentation is deliberately top-down: we begin
with the high-level shape of the pipeline and the single idea that organises
it, introduce the central data structure, the Intent Abstract Model, and
the function that consumes it, state the soundness theorem that the whole
construction is built to satisfy, and close with the principal design
decisions, each paired with the gap or objective it answers. Implementation
detail is deferred to Section~\ref{rem:impl}; here the concern is the
\emph{why} and the \emph{shape}, not the \emph{how}.

\subsection{Design rationale}
\label{sec:rationale}

One observation drives the architecture. The corrective transformation we
need (moving data from the grammar to the protocol) can only be expressed
if the tool knows which parts of a constructed command are fixed skeleton
and which are variable data. That knowledge is precisely what a detector
withholds (gap G2). It follows that the engine must \emph{reconstruct} the
developer's intent as an explicit, typed object, rather than editing the
host syntax tree directly. Everything else is downstream of this choice. If
intent is reconstructed faithfully, a sound rewrite becomes a matter of
binding each data position to a parameterising API; if intent cannot be
reconstructed faithfully, the only safe move is to abstain. The architecture
is therefore a funnel: raw code in, an intent model in the middle, and
either a certified patch or a reasoned \bottom{} out.

A second, quieter principle is the separation of \emph{producing} a patch
from \emph{accepting} one. We do not trust the component that generates an
edit, whatever it is, to also certify the edit. Acceptance instead hinges on
a decidable property of the patched code that an independent validator
re-checks. This decoupling has a pleasant consequence: the soundness of an
accepted patch does not depend on the correctness of the synthesiser, only
on the correctness of the much smaller checker. It is a classic
trust-minimising move, and it underwrites the auditability objective R6.

\subsection{The seven-stage pipeline}
\label{sec:pipeline}

\irsam{} is organised as a linear pipeline of seven stages, which we name
A through G and compose right-to-left as $G\circ F\circ E\circ D\circ C\circ
B\circ A$. Each stage has a typed interface, transforms its input into the
input expected by the next stage, and (this is the recurring
theme) may abstain by returning \bottom{} with a reason.
Figure~\ref{fig:pipeline} shows the arrangement and singles out the two
stages, D and E, that form the conceptual core. Stages A through C are
preparatory: they turn a finding and a file into a parameterised template.
Stage D lifts that template into the intent model. Stage E binds the model
to a concrete rewrite. Stages F and G realise and certify the result.

\begin{figure}[ht]
\centering
\includegraphics[width=\linewidth]{figures/fig_pipeline.pdf}
\caption{The \irsam{} seven-stage pipeline. A finding enters at stage~A and
is progressively refined into a parameterised template (A--C), an Intent
Abstract Model (D), a patch plan (E), a synthesised host-AST edit (F) and
finally a certified patch (G). The dashed region encloses the conceptual
core: intent reconstruction and the binding function $\varphi$. Any stage
may return $\bot$ with a machine-readable reason instead of passing output
downstream.}
\label{fig:pipeline}
\end{figure}

The typed signatures make the contract between stages explicit, and they are
worth stating once because the implementation in Section~\ref{rem:impl}
mirrors them line for line:
\begin{description}[leftmargin=1.7em,style=nextline]
  \item[A: Sink locator.] $\mathrm{Code}\to\mathcal{P}(\mathrm{Finding})$.
  Consumes a normalised finding and resolves it to a concrete sink call-site.
  \item[B: Slicer.]
  $\mathrm{Finding}\times\mathrm{Code}\to\mathcal{P}(\mathrm{Slice})$.
  Computes a backward slice of the string-construction feeding the sink;
  abstains if the slice escapes the local scope.
  \item[C: Reconstructor.] $\mathrm{Slice}\to\mathrm{Template}$. Rebuilds
  a parameterised template from the slice, marking data positions as holes;
  abstains on grammar outside the supported fragment.
  \item[D: Intent parser.]
  $\mathrm{Template}\times\mathrm{Interpreter}\to\mathrm{IAM}\cup\{\bottom\}$.
  Parses the template into an Intent Abstract Model, disambiguating where
  necessary.
  \item[E: Binder ($\varphi$).]
  $\mathrm{IAM}\to\mathrm{Patch}\cup\{\bottom\}$. Binds each hole to a safe
  API through the verified catalog.
  \item[F: Synthesiser.]
  $\mathrm{Patch}\times\mathrm{HostAST}\to\mathrm{HostAST}'$. Applies the
  patch plan to the host syntax tree.
  \item[G: Validator.]
  $\mathrm{HostAST}'\to\{\textsf{accept},\textsf{reject}\}$. Re-checks the
  patched program through five independent gates.
\end{description}

\subsection{Component architecture}
\label{sec:components}

The pipeline stages map onto a layered code architecture, summarised in
Figure~\ref{fig:architecture}. At the top, an input layer turns the output
of an external detector into the unified finding schema; this is where the
detector-agnostic objective R4 is realised. The core layer hosts the seven
stages as a set of cooperating modules: slicing, reconstruction, parsing,
disambiguation, the intent model, the binder and the rewriter, all driven by
a thin pipeline orchestrator. Beneath the core sits a knowledge-and-proofs
layer: the declarative binder catalogs, the JSON schema that validates them,
and the Lean~4 development that mechanises the soundness argument. At the
bottom, the interface layer exposes a command-line tool and the evaluation
harnesses. The deliberate thinness of the core, and the externalisation of
safe-API knowledge into data, are recurring choices we justify in
Section~\ref{sec:decisions}.

\begin{figure}[ht]
\centering
\includegraphics[width=\linewidth]{figures/fig_architecture.pdf}
\caption{Layered component architecture of \irsam{}. The input layer
normalises detector output into a unified finding schema; the core layer
implements the seven pipeline stages as cooperating modules under a thin
orchestrator; the knowledge-and-proofs layer externalises safe-API knowledge
into declarative catalogs validated by a JSON schema and backs the soundness
argument with Lean~4 proofs; the interface layer exposes the command-line
tool and the evaluation harnesses. Solid arrows denote data flow; dashed
arrows denote that the core consults the knowledge layer.}
\label{fig:architecture}
\end{figure}

\subsection{The Intent Abstract Model}
\label{sec:iam}

The object at the centre of Figure~\ref{fig:pipeline} is the Intent Abstract
Model, and it is the single most important idea in the design. We define it
following the formal model that the implementation realises.

\begin{definition}[Intent Abstract Model]
\label{def:iam}
An Intent Abstract Model is a 4-tuple $A=(G,H,C,\Sigma)$ in which
$G$ is a typed \emph{Semantic Intent Graph} (SIG) conforming to a schema
$S_I$ over the interpreter's grammar, containing ordinary syntax nodes and
explicit \textsc{hole} nodes; $H$ maps each hole to its grammatical role
$\in\{\textsf{value},\textsf{identifier},\textsf{fragment},
\textsf{structural}\}$, its semantic type
$\in\{\textsf{string},\textsf{integer},\textsf{decimal},\textsf{boolean},
\textsf{date},\textsf{blob},\textsf{enum}\}$ and its cardinality
$\in\{\textsf{one},\textsf{many(bounded)},\textsf{many(unbounded)}\}$;
$C$ is a finite set of constraints (type, range, equality, allow-list) over
the holes; and $\Sigma$ is the symbol environment carrying, for every
identifier occurring in $G$, its host-language definition together with the
proof obligation that the identifier resolves to a value in the trust domain
$D_T$.
\end{definition}

The Semantic Intent Graph is a typed abstract syntax tree of the
\emph{interpreter command}, not of the host language, in which every
position where host data flows in is marked as a typed hole.
Figure~\ref{fig:iam-sig} shows the SIG for a representative query with two
value holes. The crucial property is that the SIG makes the developer's
intent a first-class object: the fixed skeleton (the \code{Select},
\code{From} and \code{Where} structure) is cleanly separated from the
variable data (the two \textsf{value} holes), directly resolving the
impedance mismatch of gap G2. A repairer that holds the SIG no longer has to
guess which substring is data; the answer is written into the structure.

\begin{figure}[ht]
\centering
\includegraphics[width=0.92\linewidth]{figures/fig_iam_sig.pdf}
\caption{The Semantic Intent Graph for the query
\code{SELECT id FROM users WHERE name = $\langle$h1$\rangle$ AND role =
$\langle$h2$\rangle$}. Blue nodes are interpreter-grammar syntax; white
nodes are developer-controlled literals and identifiers; green nodes are
typed holes that mark the positions where host data is inserted. Hole
\code{h1} is a \textsf{value} hole of semantic type \textsf{string}; hole
\code{h2} is a \textsf{value} hole of semantic type \textsf{enum}. The
fixed skeleton and the variable data are thereby separated structurally.}
\label{fig:iam-sig}
\end{figure}

The four-way classification of holes carries the weight of the soundness
argument, so it merits a closer look; Table~\ref{tab:holes} lays it out.
The decisive line is the first: a \textsf{value} hole is the only kind that
may legitimately be bound to a runtime expression drawn from $D_A$, because
a value can be relocated to the protocol channel and rendered inert. Every
other kind (an identifier such as a table or column name, or a structural
fragment such as an ordering direction) cannot be parameterised by most
drivers and must instead be proven to come from $D_T$, typically via an
allow-list. When that proof is unavailable, the only sound response is to
abstain. This table is, in a sense, the whole policy of the engine
compressed into four rows.

\begin{table}[ht]
\centering
\caption{The four grammatical roles a hole may take, and the sound
realisation strategy for each. Only \textsf{value} holes may be bound to
attacker-controlled data; the rest require a trust proof or force
abstention.}
\label{tab:holes}
\begin{tabular}{p{2.4cm}p{4.7cm}p{5.0cm}}
\hline
\textbf{Role} & \textbf{Example position} & \textbf{Sound realisation} \\
\hline
\textsf{value} & \code{WHERE name = $\langle$h$\rangle$} &
Bind on the parameter channel (e.g.\ \code{setString}); may carry $D_A$. \\
\textsf{identifier} & \code{ORDER BY $\langle$h$\rangle$} (column) &
Allow-list lookup over $D_T$; trap on miss. \\
\textsf{fragment} & a reusable sub-clause &
Bind only if every leaf is itself soundly realisable. \\
\textsf{structural} & \code{ASC}/\code{DESC}, set operators &
Map to a fixed finite vocabulary in $D_T$; otherwise abstain. \\
\hline
\end{tabular}
\end{table}

\subsection{The binding function \texorpdfstring{$\varphi$}{phi}}
\label{sec:phi}

If the SIG is the noun of the design, the binding function is its verb.

\begin{definition}[Binding function $\varphi$]
\label{def:phi}
$\varphi:\mathrm{IAM}\rightharpoonup\mathrm{Patch}$ is a partial function
that maps an Intent Abstract Model to a host-language rewrite through a
finite catalog of verified \emph{binders}. Each binder is a tuple
$(\mathit{Pattern},\mathit{Rewrite},\mathit{Pre},\mathit{Post},\mathit{Proof})$,
pairing a SIG sub-pattern with a safe-API rewrite template, checkable
pre- and post-conditions over the host AST, and a soundness argument for
the rewrite. By specification, $\varphi(A)=\bottom$ whenever any hole of
$A$ cannot be covered by a sound binder.
\end{definition}

Operationally, $\varphi$ matches maximal sub-graphs of the SIG against the
catalog with a deterministic, greedy, left-to-right tie-break and composes
the resulting AST rewrites, provided every hole is covered. The partiality
is the point. An identifier hole whose trust cannot be established, or a
construction outside the supported grammar fragment, leaves a hole
uncovered and forces \bottom{}. Abstention here is not an exceptional error
path bolted onto the side; it is a normal, specified output of the function,
which is what realises objective R2. One may reasonably ask why the catalog
is finite and externalised rather than learned. The answer is auditability:
a finite catalog of human-reviewed binders, each with a stated proof
obligation, can be checked; a learned mapping cannot.

\subsection{The decidable soundness criterion}
\label{sec:thm}

The architecture culminates in a theorem. It connects the structural shape
of a patch (something a machine can inspect) to the semantic guarantee of
Eq.~\eqref{eq:soundness}.

\begin{theorem}[Structural soundness]
\label{thm:soundness}
Let $p=\varphi(A)$ be a patch. If, in $p$, every \textsf{value} hole of $A$
is realised as an argument to a parameterising API call of the sink, and
every \textsf{identifier}/\textsf{structural} hole is realised as an
allow-list lookup over $D_T$ that traps on miss, then $p$ satisfies
Eq.~\eqref{eq:soundness}: no execution of $p$, on any input, lexes
attacker-controlled bytes as interpreter syntax.
\end{theorem}

The significance of Theorem~\ref{thm:soundness} lies less in the statement
than in the nature of its hypothesis. That hypothesis is a \emph{decidable}
property of the patched host AST: one can check, mechanically and without
running the program, whether each value hole became a parameter argument and
each identifier hole became a guarded lookup. Consequently the validator of
stage~G can re-certify any patch independently of how it was produced,
which is exactly the trust-minimising decoupling promised in
Section~\ref{sec:rationale} and the substance of objective R6. The theorem
is proved on paper for the $\mathrm{SQL}_0$+JDBC fragment and mechanised in
Lean~4 for the identifier-allowlist closure, the structural cover, the audit
soundness and the principal parameterisation result; the OS-command
argument-inertness results and the path-guard insertion lemmas are likewise
mechanised. Section~\ref{sec:lean} returns to the mechanisation in detail.
The argument rests on one supporting lemma, which isolates the single
external assumption of Section~\ref{sec:threat}.

\begin{lemma}[Parameter inertness, JDBC]
\label{lem:inertness}
A value bound through \code{PreparedStatement.setX} is transmitted on the
JDBC parameter channel and is never re-parsed as SQL by the driver, under
the dialect-prepared-statement assumption. Hence a \textsf{value} hole
realised in this way cannot contribute syntax to the executed command.
\end{lemma}

A short remark on scope is in order. The fragment for which the paper proof
is complete, $\mathrm{SQL}_0$, is a restricted but practically common subset
of SQL: single-statement \code{SELECT}/\code{INSERT}/\code{UPDATE}/\code{DELETE}
with equality, \code{IN}, \code{LIKE}, \code{LIMIT} and \code{ORDER BY}.
Commands outside this subset are not handled unsoundly; they are simply not
handled, and the engine abstains. This is a conscious trade of breadth for
certainty, and Section~\ref{rem:discussion} weighs it honestly.

\subsection{Design decisions}
\label{sec:decisions}

We close the architectural account by recording the principal design
decisions and the rationale behind each. They are numbered D1 through D6,
and each points back to a gap from Section~\ref{sec:gap} or an objective
from Section~\ref{rem:intro}. Listing them explicitly is not bookkeeping for
its own sake; it makes the chain from problem to mechanism inspectable,
which is itself a form of the auditability we advocate.

\paragraph{D1: Reconstruct intent rather than rewrite bytes.}
We lift the command into a typed SIG instead of editing the host AST
directly. Byte- or AST-level rewriting cannot tell skeleton from data on
dynamically assembled commands (G2); intent reconstruction makes the
distinction explicit and is what enables a provable transformation in the
first place. This is the foundational decision; the others elaborate it.

\paragraph{D2: Make abstention first-class.}
Every stage may return \bottom{} with a reason. This follows directly from
the anti-false-fix premise (G3): a pipeline that must always emit an edit
cannot be sound. Because the reasons are machine-readable, a human or a
stronger tool can resume precisely where \irsam{} stops, rather than
re-triaging from scratch.

\paragraph{D3: Accept by a decidable, re-checkable criterion.}
Rather than trust whatever produced a patch, acceptance hinges on the
hypothesis of Theorem~\ref{thm:soundness}, which the validator re-checks on
the patched AST. This decouples how a patch is generated from whether it is
accepted and yields the auditable certificate of R6. It also means the
synthesiser can be improved, or even replaced, without weakening any
guarantee.

\paragraph{D4: A declarative binder catalog.}
Safe-API knowledge is externalised into YAML catalogs, one per
framework-interpreter pair, rather than hard-coded. Each binder declares its
SIG pattern, rewrite template, parameterising APIs and proof obligation.
Extending coverage then becomes a data change validated against a schema,
not a code change, which supports R5 and keeps the verified core small. A
small core is easier to trust; that is the entire motivation.

\paragraph{D5: Detector-agnostic input.}
Stage~A consumes a unified finding schema normalised from SARIF
\cite{sarif2023}. Detection quality and remediation soundness thereby become
independent concerns (R4): a weak bundled detector limits how \emph{many}
findings \irsam{} acts on, never \emph{whether} its accepted patches are
sound. We regard this separation as essential; conflating the two is a
common source of confusion in the literature.

\paragraph{D6: Independent multi-gate validation.}
Acceptance requires passing five orthogonal gates: compile, regression,
structural, re-SAST and differential. The redundancy is intentional. The
structural gate certifies Theorem~\ref{thm:soundness}; the others guard
against accidental breakage and confirm that the original finding no longer
fires. No single gate is load-bearing for soundness, but together they make
an accepted patch trustworthy along several independent axes.


% =====================================================================
\section{Implementation}
\label{rem:impl}
% =====================================================================

We now move from architecture to artefact. This section walks through the
implementation of \irsam{} in enough depth to make the design concrete and
the evaluation reproducible. The order roughly follows the flow of data: we
start with the technology stack and the reasons behind it, describe the
unified finding schema that admits any detector, traverse the seven stages
as they appear in the codebase, examine the declarative binder catalogs and
their schema, detail the five validation gates, summarise the Lean~4
mechanisation, and end with the command-line surface and a fully worked
transformation. Throughout, we name the actual modules and data structures
so that the prose can be checked against the source.

\subsection{Technology stack and choices}
\label{sec:stack}

\irsam{} is implemented in Python~3.10. The choice is pragmatic rather than
ideological: the project must parse and manipulate both Java and Python
source, orchestrate external tools (SAST detectors, the Java compiler, a
database used for differential testing), and present a usable command-line
interface. Python's ecosystem covers all of these with little friction, and
its dynamic typing was tempered in practice by Pydantic data classes and a
\code{mypy} type-checking gate. Table~\ref{tab:stack} summarises the stack
and the role of each component.

\begin{table}[ht]
\centering
\caption{Implementation and evaluation stack of \irsam{}.}
\label{tab:stack}
\begin{tabular}{p{4.6cm}p{8.6cm}}
\hline
\textbf{Concern} & \textbf{Choice} \\
\hline
Core language          & Python 3.10 \\
SQL$_0$ grammar parser & \code{lark} (Earley/LALR) \\
Shell / XPath parsing  & \code{antlr4} runtime \\
Data validation        & \code{pydantic} (typed records) \\
CLI framework          & \code{typer} (\code{scan/plan/patch/validate/pipeline}) \\
Terminal reporting     & \code{rich} tables \\
Binder catalogs        & declarative YAML, JSON-schema validated \\
Detector (default)     & Semgrep 1.164.0, offline bundled ruleset \cite{semgrep2024} \\
Java compile gate      & JDK 21 \code{javac} (real compilation) \\
Python compile gate    & in-process CPython \code{compile} \\
Differential gate DB   & SQLite (in-memory) \\
Formal proofs          & Lean 4 \cite{demoura2021lean4} with mathlib \cite{mathlib2020} \\
Test framework         & \code{pytest} (unit, property via \code{hypothesis}, integration) \\
\hline
\end{tabular}
\end{table}

Two choices deserve a word of justification, since they shape the
trustworthiness of the result. Encoding safe-API knowledge in YAML, and
validating each catalog against \code{schemas/binder.schema.json}, keeps the
logic that must be trusted small and turns coverage extension into a data
change. Eleven catalogs ship at the time of writing, spanning JDBC, Spring
\code{JdbcTemplate}, MyBatis, Hibernate HQL, JPA, Django ORM, Python DB-API,
Python \code{subprocess}, Java \code{ProcessBuilder}, LDAP and XPath.
Separately, the decision to mechanise the soundness argument in Lean~4 was
made because a security claim of this kind is only as credible as its proof;
a dependently typed kernel lets us discharge the central lemmas with no
\code{sorry} placeholders, a point we substantiate in
Section~\ref{sec:lean}.

\subsection{The unified finding schema}
\label{sec:schema}

Stage~A's job is to make the rest of the pipeline indifferent to which
detector raised an alert. It normalises a SARIF finding
\cite{sarif2023} into an internal record (a Pydantic model) that carries the
sink location (file and line), the host language, the interpreter
(\code{sql}, \code{shell}, \code{ldap} or \code{xpath}) and the sink API. A
rule-map translates detector-specific rule identifiers, including the
namespaced identifiers that Semgrep assigns to local-file rules, into this
record. The effect is that a CodeQL finding \cite{avgustinov2016codeql}, a
Semgrep finding \cite{semgrep2024} and a SonarQube finding all converge on
the same downstream representation, after which they are indistinguishable
to the engine. This is the concrete realisation of design decision~D5, and
it is the mechanism behind the detector-agnostic objective R4. The
\code{ingest/} module owns this translation; everything below it sees only
the normalised record.

\subsection{Stage-by-stage realisation}
\label{sec:stages}

We trace the seven stages as they are implemented, naming the module that
owns each. The reader may find it helpful to keep Figure~\ref{fig:pipeline}
in view, since the code structure deliberately mirrors the typed signatures
given in Section~\ref{sec:pipeline}.

\subsubsection{Stage~B: Slicing}
The slicer (\code{core/slicer/}) computes, for the language at hand, a
backward slice of the expression feeding the sink argument. Its central
routine, \code{slice\_sink\_argument}, returns a \code{SliceResult} carrying
the reconstructed skeleton, the sink call text and its line. Where a helper
method assembles part of the query, the slicer performs a single hop of
inter-procedural inlining; beyond one hop, or when the construction depends
on state outside the local method, it returns a \code{SliceAbstention} with
a reason such as \code{no connection variable in scope} or \code{slice
escapes scope}. The reason this matters is subtle: reconstructing intent
from a slice that leaves the local scope would require assumptions the
engine cannot verify, so abstaining is the sound choice. As
Section~\ref{rem:results} shows, this particular abstention is the single
most common outcome on real code, which is itself an informative finding.

\subsubsection{Stage~C: Reconstruction}
The reconstructor (\code{core/recon/}) rebuilds a parameterised template
from the slice, producing a \code{ParameterizedTemplate}. It recognises
string concatenation and, for Python, f-strings as data-insertion points,
turning each into a hole position. Constructions outside this grammar
fragment (Python \code{\%}-formatting against a dictionary, say, or SQL
\code{CallableStatement} \code{\{call \dots\}} syntax) trigger a
\code{ReconAbstention} rather than a guessed reconstruction. The guiding
maxim is the one from Section~\ref{sec:premise}: a precise abstention beats
an imprecise repair.

\subsubsection{Stage~D: Intent parsing and disambiguation}
The intent parser (\code{core/parsers/}) lifts the template into a SIG over
the relevant interpreter grammar ($\mathrm{SQL}_0$ for SQL, with analogous
schemas for shell, LDAP and XPath), yielding an \code{IAM} or raising a
typed error such as \code{SQL0SyntaxError} or \code{SQL0AmbiguousIntent}.
When a token is genuinely ambiguous, an optional disambiguator
(\code{core/disambig/}) selects one label from a finite, syntactically safe
set and re-invokes the parser. The disambiguator is bounded by
construction: it can only choose among labels that are individually sound,
so even a wrong guess cannot produce an unsound parse. If parsing still
fails, the stage abstains, annotating the original reason with the
disambiguator's decision for provenance. We regard this ``safe-by-design''
framing of disambiguation as one of the quieter but more important
engineering ideas in the system.

\subsubsection{Stage~E: Binding ($\varphi$)}
The binder (\code{core/phi/} together with \code{core/binder/}) loads the
catalog appropriate to the $(\text{language},\text{interpreter})$ pair,
matches SIG sub-patterns against the binders, and produces a
\code{PatchPlan}. The plan is a structured object, not a string: it records
the prepared template, a tuple of \code{SetterCall}s (each naming a
parameter index, an API, the host expression and the originating hole), any
\code{AllowlistGuard}s for identifier holes, the per-hole
\code{HostRealization}s, and the proof obligations the binders carry. A
value hole, for example, becomes a \code{setString(i, expr)} against a
\code{PreparedStatement}. Should any hole remain uncovered, $\varphi$
returns a \code{PhiAbstention}. The structured plan is what later lets the
validator re-check soundness without re-deriving it.

\subsubsection{Stage~F: Synthesis}
The synthesiser (\code{core/rewriter/}) applies the plan to the host source,
emitting the patched text and a unified diff inside a \code{PatchResult}.
For Java it discovers the connection variable in scope, introduces a
\code{PreparedStatement}, threads the \code{setX} calls, and removes the now
dead query-string variable; for Python it rewrites the cursor call to a
parameterised form with \code{\%s} placeholders. Backends are registered for
the full A--G path only for the $(\text{java},\text{sql})$ and
$(\text{python},\text{sql})$ pairs. The shell, LDAP and XPath interpreters
are parsed and modelled (they have SIG schemas and binders) but no
registered rewriter yet, so they yield the typed \code{unsupported\_backend}
abstention. We stress that this is a sound refusal, not a crash: the engine
knows what it would need to do and declines because it cannot yet do it
provably.

\subsubsection{Stage~G: The five gates}
\label{sec:gates}
A patch is accepted only if it passes all five gates of
\code{core/validator/}, illustrated in Figure~\ref{fig:gates}. We describe
each in turn, since their independence is the substance of decision~D6.
\begin{enumerate}[leftmargin=1.9em,label=\arabic*.]
  \item \textbf{Compile.} For Java the gate invokes \code{javac} in a
  temporary directory and performs a real compilation; for Python it uses
  the in-process CPython \code{compile} built-in. An early defect that ran
  \code{javac} unconditionally, and therefore failed every Python
  patch, was found and fixed during evaluation (Section~\ref{sec:iteration}).
  \item \textbf{Regression.} When the project ships a test runner (Maven,
  Gradle, \code{pytest} or npm, detected by manifest), the gate runs it
  within an allow-list of permitted commands and a timeout, and the existing
  tests must still pass. The allow-list is a small security boundary: the
  engine never executes an arbitrary command from an untrusted repository.
  \item \textbf{Structural.} The gate re-checks the hypothesis of
  Theorem~\ref{thm:soundness} directly on the patched AST: every value hole
  realised through a parameterising API in the trusted set, every identifier
  hole through an allow-list lookup, every symbol proven in $D_T$. This gate
  \emph{is} the soundness certificate.
  \item \textbf{Re-SAST.} The detector is re-run, and the gate confirms that
  the original finding class no longer fires, with no residual concatenation in
  the sink argument.
  \item \textbf{Differential.} The original and patched statements are
  executed against an in-memory SQLite fixture over both benign and attack
  payloads. Benign inputs must yield identical row-sets; attack payloads
  must neither execute additional statements nor mutate the schema. A second
  defect, which mislabelled an inconclusive baseline as a regression, was
  found and fixed here as well (Section~\ref{sec:iteration}).
\end{enumerate}

\begin{figure}[ht]
\centering
\includegraphics[width=\linewidth]{figures/fig_gates.pdf}
\caption{The five-gate validator of stage~G. A synthesised patch must pass
every gate in sequence to be accepted; failing any one gate causes
rejection, never a downgraded acceptance. The structural gate (highlighted)
certifies the hypothesis of Theorem~\ref{thm:soundness} on the patched AST,
while the compile, regression, re-SAST and differential gates guard against
accidental breakage and confirm that the original finding no longer fires.}
\label{fig:gates}
\end{figure}

Figure~\ref{fig:sequence} draws the same process as a sequence of messages,
which makes the abstention behaviour vivid: at every arrow, a stage that
cannot discharge its proof obligation returns \bottom{} rather than passing
output along. A repair is thus a relay in which any runner may, legitimately,
decline to continue.

\begin{figure}[ht]
\centering
\includegraphics[width=\linewidth]{figures/fig_sequence.pdf}
\caption{Sequence view of a single repair. The detector's finding is refined
into a parameterised template, an Intent Abstract Model, a patch plan and a
patched AST, and is finally subjected to the five gates. At any message
boundary a stage may return $\bot$ with a machine-readable reason instead of
forwarding output, so that the overall outcome is either an accepted,
certified patch or a reasoned abstention.}
\label{fig:sequence}
\end{figure}

\subsection{Key data structures}
\label{sec:datastructures}

It is worth making the central data structures explicit, because their field
names recur in the evaluation. The Intent Abstract Model is a frozen record
holding the SIG, a tuple of holes, the constraints and the symbol
environment, tagged with the interpreter. Each hole carries its name,
grammatical context, semantic type, cardinality and, for identifier
holes, an optional allow-list. The symbol environment maps each identifier
to its host expression and a boolean recording whether taint analysis placed
it in $D_T$. The binder's output, the patch plan, is likewise a frozen
record, and its proof-obligation field is what the structural gate consumes.
The following sketch, lightly abbreviated from \code{core/iam/} and
\code{core/phi/}, conveys the shape without the full type annotations:

\begin{verbatim}
IAM        = (sig, holes, constraints, symbols, interpreter)
Hole       = (name, ctx, sem, card, allowlist)
              ctx  in {value, identifier, fragment, structural}
              sem  in {string, integer, decimal, boolean, date, blob, enum}
              card in {one, many.bounded, many.unbounded}
SymbolEntry= (name, host_expr, proven_in_DT, constant_value)
PatchPlan  = (prepared_template, setter_calls, allowlist_guards,
              realizations, catalog_id, binder_ids_used, proof_obligations)
\end{verbatim}

The audit logic that the structural gate invokes checks four predicates,
P1 through P4: syntactic well-formedness with no residual string
concatenation; semantic-type preservation between a hole and its setter;
hole coverage, meaning every hole is realised through a parameterising API
or an allow-list; and interpreter equivalence on benign payloads, which the
differential gate establishes. The audit returns counterexamples when a
predicate fails, which proved valuable when diagnosing the gate defects of
Section~\ref{sec:iteration}.

\subsection{The declarative binder catalogs}
\label{sec:catalogs}

The eleven binder catalogs are the externalised knowledge of the system, and
Table~\ref{tab:binders} inventories them. Each catalog declares the
interpreter and host language it applies to, the framework it targets and
the parameterising APIs it is allowed to emit. The loader
(\code{core/binder/loader.py}) validates every catalog against the JSON
schema and then performs a closure check: every API referenced in a rewrite
must appear in the catalog's declared \code{parameterizing\_apis}. This
check is what prevents a careless catalog edit from introducing an API the
soundness argument does not cover.

\begin{table}[ht]
\centering
\caption{The eleven binder catalogs shipped with \irsam{}. Catalogs marked
``modelled'' have SIG schemas and binders but no registered host rewriter
yet, so they yield sound \code{unsupported\_backend} abstentions.}
\label{tab:binders}
\begin{tabular}{p{3.6cm}p{1.6cm}p{1.6cm}p{4.6cm}}
\hline
\textbf{Catalog} & \textbf{Interp.} & \textbf{Lang.} & \textbf{Parameterising API} \\
\hline
\code{sql\_jdbc}            & SQL   & Java   & \code{PreparedStatement.setX} \\
\code{spring\_jdbctemplate} & SQL   & Java   & \code{JdbcTemplate} args \\
\code{mybatis}              & SQL   & Java   & mapper parameter binding \\
\code{hibernate\_hql}       & SQL   & Java   & \code{Query.setParameter} \\
\code{jpa}                  & SQL   & Java   & \code{TypedQuery.setParameter} \\
\code{sql\_pydbapi}         & SQL   & Python & \code{cursor.execute} (\code{\%s}) \\
\code{django\_orm}          & SQL   & Python & QuerySet parameterisation \\
\code{python\_subprocess}   & shell & Python & \code{subprocess} argv, \code{shell=False} \\
\code{java\_processbuilder} & shell & Java   & \code{ProcessBuilder(List)} \\
\code{ldap} (modelled)      & LDAP  & any    & RFC-4515 escaping helpers \\
\code{xpath} (modelled)     & XPath & any    & variable-resolver APIs \\
\hline
\end{tabular}
\end{table}

A single binder, in YAML, makes the structure concrete. The following
excerpt from \code{sql\_jdbc.yaml} binds an equality predicate whose
right-hand side is a value hole, rewriting it to a placeholder and choosing
the setter by the hole's semantic type; its proof obligation references
Lemma~\ref{lem:inertness}:

\begin{verbatim}
- id: sql-eq-value
  pattern:
    kind: Comparison
    children:
      - {kind: Column, bind: col}
      - {kind: Token, lit: "="}
      - {kind: Hole, ctx: value,
         sem_in: [string, integer, decimal, boolean, date], bind: v}
  rewrite:
    template_string: "{col} = ?"
    bindings:
      - {param_index: $auto, host_expr: "$v.host_expr",
         api: "choose_setter($v.sem)"}
  postconditions:
    - structural_value_via_parameterized_api: $v
  proof_obligation: "Lemma 1 (Parameter inertness)."
\end{verbatim}

The pattern is a SIG sub-graph schema with named binds; the rewrite is a
host-AST template referencing those binds; the postcondition is exactly the
structural property the gate will later verify. Adding a new safe API is, by
design, a matter of authoring such a record and letting the schema and
closure checks vouch for its well-formedness.

\subsection{Mechanised soundness in Lean~4}
\label{sec:lean}

A security guarantee is only as strong as its proof, so the soundness
argument is mechanised in Lean~4 \cite{demoura2021lean4}, with mathlib as a
foundation \cite{mathlib2020}. The development lives under \code{formal/IRSAM/}
and is organised into a string-algebra core, per-interpreter models for SQL
and shell, binder models for JDBC and \code{subprocess}, and a soundness
layer whose modules mirror the paper proof. Figure~\ref{fig:lean} summarises
the status. Twenty theorems and lemmas are closed with no \code{sorry} in
the source tree, including the identifier-allowlist closure, the
structural cover, the audit soundness and completeness results, the shell
argument-inertness family, and the two principal results
\code{soundness\_parameterize} and \code{ir\_sam\_soundness}. A further eight
statements have landed with placeholder \code{trivial} proofs; these are the
per-framework parameter-inertness lemmas (Hibernate, JPA, MyBatis, Spring,
and the JDBC and DB-API wire-protocol models) whose full discharge requires
an equational model of the driver protocol and is left as future work. The
acceptance gate \code{scripts/check\_no\_sorry.sh} confirms that no
\code{sorry} appears in the proved core, and the kernel reports only the
three standard classical axioms (\code{propext}, \code{Classical.choice} and
\code{Quot.sound}).

\begin{figure}[ht]
\centering
\includegraphics[width=\linewidth]{figures/fig_lean_status.pdf}
\caption{Status of the Lean~4 mechanisation. \emph{Left:} of the twenty-eight
statements in the development, twenty are fully closed with no \code{sorry},
eight are partial (statement landed, proof \code{trivial}), and none are
deferred via \code{sorry}. \emph{Right:} per-module breakdown of closed and
partial results. The principal parameterisation and audit-soundness theorems
are among the closed set; the open eight are per-framework wire-protocol
inertness lemmas.}
\label{fig:lean}
\end{figure}

We are candid that the eight partial lemmas are a genuine gap rather than a
formality: until the JDBC wire-protocol model is supplied, the mechanised
guarantee for those frameworks rests on the paper argument plus the closed
structural results, not on an end-to-end Lean proof. What \emph{is}
mechanised, however, is the part the validator relies on, namely that a patch
satisfying the structural hypothesis is sound, and so the acceptance criterion
itself stands on proved ground.

\subsection{Command-line surface}
\label{sec:cli}

The engine is driven through a Typer command-line interface
(\code{core/cli.py}) exposing five subcommands that correspond to prefixes of
the pipeline, listed in Table~\ref{tab:cli}. The progressive design is
convenient in practice: an engineer can inspect what the slicer and binder
infer with \code{plan} before ever writing a file, and can run the full
gated pipeline with \code{pipeline} when a certified patch is wanted. The
same entry points back the evaluation harnesses.

\begin{table}[ht]
\centering
\caption{The five command-line subcommands and the pipeline stages each one
exercises.}
\label{tab:cli}
\begin{tabular}{p{2.2cm}p{2.2cm}p{8.0cm}}
\hline
\textbf{Command} & \textbf{Stages} & \textbf{Effect} \\
\hline
\code{scan}     & A     & List the sink call-sites detected in a file. \\
\code{plan}     & A--E  & Produce and print a patch plan without rewriting. \\
\code{patch}    & A--F  & Write the patched source and an optional diff. \\
\code{validate} & A--G  & Run the five gates on a patched file. \\
\code{pipeline} & A--G  & Full end-to-end run, emitting a JSON report. \\
\hline
\end{tabular}
\end{table}

\subsection{A worked transformation}
\label{sec:worked}

To ground the whole pipeline in a single concrete example, consider the
canonical CWE-89 sink below, in which a value is concatenated into the SQL
grammar. The transformation \irsam{} performs relocates that value to a
protocol-bound parameter slot, exactly the move depicted abstractly in
Figure~\ref{fig:dataflow}:

\begin{verbatim}
// BEFORE  (value concatenated into the SQL grammar)
String sql = "SELECT id FROM users WHERE name = '" + name + "'";
ResultSet rs = stmt.executeQuery(sql);

// AFTER   (value relocated to a protocol-bound parameter slot)
PreparedStatement ps =
    conn.prepareStatement("SELECT id FROM users WHERE name = ?");
ps.setString(1, name);
ResultSet rs = ps.executeQuery();
\end{verbatim}

Tracing this through the stages: the slicer recovers the concatenated
construction; the reconstructor marks \code{name} as a hole and produces the
template \code{SELECT id FROM users WHERE name = ?}; the intent parser lifts
it into the SIG of Figure~\ref{fig:iam-sig} (with one value hole); the
binder matches the \code{sql-eq-value} binder above and plans a
\code{setString} call; the synthesiser emits the rewritten Java; and the
validator certifies that the lone value hole is realised as a parameter
argument (the hypothesis of Theorem~\ref{thm:soundness}), so the structural
gate accepts. The example is small, but every stage in
Figure~\ref{fig:pipeline} participates, and the resulting patch carries a
certificate a reviewer can re-check by inspection.


% =====================================================================
\section{Experimental Setup}
\label{rem:setup}
% =====================================================================

The evaluation is designed to answer the three hypotheses of
Section~\ref{sec:hypotheses} and, through them, to test whether the six
objectives are met in practice rather than merely on paper. This section
fixes the experimental conditions: the hardware and software environment,
the corpora and testbeds, the metrics, the baselines and the rationale for
choosing them, and the procedure each harness follows. A guiding principle,
worth stating up front, is that every reported figure is produced by
executing the real pipeline end to end: compiling with a real \code{javac},
scanning with a real Semgrep, and exercising a real SQLite oracle. None of
the numbers below is hand-authored or simulated.

\subsection{Environment}
\label{sec:environment}

All experiments were run on a single Windows workstation with Python~3.10,
JDK~21 (so the compile gate performs genuine Java compilation rather than a
syntactic approximation), Semgrep~1.164.0 with its offline bundled ruleset
\cite{semgrep2024}, and SQLite serving as the differential fixture. Running
offline is deliberate: it removes any dependence on network conditions or
remote rule updates and makes the detection figures reproducible, at the
cost of using a deliberately conservative ruleset. The unit and property
suite, comprising 336 tests exercised through \code{pytest} and
\code{hypothesis}, passes in full under this configuration, which we treat
as a precondition for trusting any larger-scale result.

\subsection{Corpora and testbeds}
\label{sec:corpora}

Evaluation proceeds over two tiers, complemented by a synthetic corpus for
the baseline comparison. The two-tier design is intentional: the controlled
tier exercises every supported interpreter against exact ground truth, while
the project tier subjects the engine to unmodified, real-world code where
ground truth is noisier but realism is high. Table~\ref{tab:corpora}
summarises the testbeds.

\begin{table}[ht]
\centering
\caption{Evaluation corpora and testbeds. The snippet and synthetic tiers
provide exact ground truth; the project tier provides realism on unmodified
upstream code.}
\label{tab:corpora}
\begin{tabular}{p{3.4cm}p{1.7cm}p{7.6cm}}
\hline
\textbf{Corpus} & \textbf{Size} & \textbf{Description} \\
\hline
Snippet tier        & 14 records & All four CWE classes, Java and Python,
three negative controls; lineage in NIST Juliet \cite{nistjuliet} and OWASP
Benchmark \cite{owaspbenchmark}. \\
Project tier        & 4 repos & OWASP BenchmarkJava (849 labelled cases),
java-sec-code, dvpwa, WebGoat; shallow-cloned and scanned unmodified. \\
Synthetic tier      & 20 cases & Generated SQL/LDAP/XPath cases across Java,
Python and JS/TS for the baseline comparison. \\
Bench generators    & N/A & Juliet-mini factory and loaders for CVEfixes
\cite{bhandari2021cvefixes} and Vul4J \cite{bui2022vul4j}. \\
\hline
\end{tabular}
\end{table}

The snippet corpus, \code{test-phase/snippet-tests.txt}, holds fourteen
records spanning all four CWE classes and both languages, with three
negative controls included specifically to test that the engine does not
manufacture findings or patches where none are warranted. Each record is
annotated with its CWE, language, interpreter, a vulnerability flag and an
expected outcome (\code{patched}, \code{abstain} or \code{clean}). The
project manifest, \code{test-phase/project-tests.txt}, lists intentionally
vulnerable repositories together with their in-scope CWEs and, where
available, the ground-truth label source. For the OWASP Benchmark the
shipped \code{expectedresults} CSV provides per-case labels, which lets us
compute genuine detection precision and recall rather than estimating them.

\subsection{Metrics}
\label{sec:metrics}

Detection is scored as binary classification against ground-truth labels,
yielding a confusion matrix $(\mathrm{TP},\mathrm{FP},\mathrm{TN},\mathrm{FN})$
and the usual precision $\mathrm{P}$, recall $\mathrm{R}$ and
$\mathrm{F1}=2\mathrm{PR}/(\mathrm{P}+\mathrm{R})$. Remediation is scored
with six metrics, M1 through M6, defined in Table~\ref{tab:metrics}. For the
controlled tier we additionally report two derived figures that speak
directly to the hypotheses: \emph{patch correctness}, the fraction of
repairable cases whose patch passes all five gates, and \emph{abstention
soundness}, the fraction of not-yet-supported sinks correctly abstained with
no unsound patch emitted. The latter two are, in effect, the empirical
counterparts of H1 and H2.

\begin{table}[ht]
\centering
\caption{The six remediation metrics. M2 is the headline ``did it actually
fix it'' figure; M4 and M6 jointly capture the anti-false-fix commitment.}
\label{tab:metrics}
\begin{tabular}{p{1.2cm}p{4.4cm}p{7.2cm}}
\hline
\textbf{Metric} & \textbf{Name} & \textbf{Definition} \\
\hline
M1 & Patch applicability      & The patch applies and compiles. \\
M2 & Functional fix rate      & M1 and the sink is removed while benign
behaviour is preserved. \\
M3 & Semantic equivalence     & Patched and original agree on benign inputs. \\
M4 & Residual-CWE density     & Sinks remaining in accepted patches (lower
is better). \\
M5 & Median latency           & Wall-clock median over the corpus. \\
M6 & Abstention precision     & Fraction of abstentions that were genuinely
necessary. \\
\hline
\end{tabular}
\end{table}

\subsection{Baselines}
\label{sec:baselines}

To situate \irsam{} we compare against six representative remediation
approaches, chosen to span the methodological landscape rather than to
cherry-pick weak opponents. Two are pattern- or analysis-driven autofixers:
Semgrep's autofix and a CodeQL-plus-GPT-4 pipeline in the style of recent
AI-assisted remediation \cite{githubautofix2024}. Three are learning-based
program repair systems: VulRepair \cite{fu2022vulrepair}, SeqTrans
\cite{chi2023seqtrans} and ChatRepair \cite{xia2023chatrepair}. The last is
a zero-shot large-language-model prompt, representing the simplest
generative approach \cite{placeholder-llm-repair}. The selection is
deliberate: each baseline embodies a different answer to the question ``how
should an automated tool decide what edit to make,'' and contrasting them
with \irsam{}'s prove-or-abstain stance is the point of the comparison. We
note, as a matter of fairness, that several of these systems target a
broader vulnerability class than injection; the comparison is therefore
scoped to the injection cases on which all tools are applicable.

\subsection{Procedure}
\label{sec:procedure}

Two harnesses drive the pipeline. The snippet harness
(\code{scripts/run\_snippet\_tests.py}) materialises each record to a
temporary file, runs the Semgrep stage-A detector to obtain a detection
signal scored against the label, runs the full A--G pipeline via
\code{run\_finding} so that dispatch follows the finding's
$(\text{language},\text{interpreter})$ pair, and re-scans the patched source
for residual sinks. The project harness
(\code{scripts/run\_project\_tests.py}) shallow-clones each repository, runs
Semgrep once over the in-scope source tree, filters to the in-scope CWEs,
maps OWASP Benchmark findings to their labelled cases through the shipped
CSV to compute detection metrics, and runs the pipeline on every in-scope
finding. Both harnesses emit machine-readable JSON alongside a human-readable
summary, and it is these emitted artefacts that the tables in the next
section report.


% =====================================================================
\section{Results and Analysis}
\label{rem:results}
% =====================================================================

We now report what the pipeline actually does. The presentation moves from
the controlled tier, where ground truth is exact and every supported
interpreter is represented, to the realistic project tier, where the engine
meets unmodified upstream code; it then places \irsam{} against the
baselines, recounts the defects the evaluation itself surfaced, and closes
by reading the results back against the three hypotheses. The reader should
keep one distinction in mind throughout: detection figures characterise an
\emph{interchangeable} bundled component, whereas remediation figures
characterise the engine that is the actual contribution.

\subsection{Controlled tier}
\label{sec:res-snippet}

On the snippet corpus the results are clean enough to state plainly.
Detection is perfect: no false positives on the three negative controls and
no false negatives on the eleven vulnerable snippets, for a confusion matrix
of $11/0/3/0$ and $\mathrm{P}=\mathrm{R}=\mathrm{F1}=1.00$. All four
repairable SQL sinks are rewritten and pass all five gates, giving a patch
correctness of $4/4$. The seven sinks whose interpreter currently lacks a
registered rewriter are soundly abstained, $7/7$, and no accepted patch
leaves a residual sink. Median latency sits below two milliseconds.
Table~\ref{tab:snippet} collects these figures, and Figure~\ref{fig:summary}
renders the headline guarantees.

\begin{table}[ht]
\centering
\caption{Controlled (snippet) tier results, produced by real
\code{javac}+Semgrep execution over the 14-record corpus.}
\label{tab:snippet}
\begin{tabular}{p{7.2cm}p{4.6cm}}
\hline
\textbf{Metric} & \textbf{Value} \\
\hline
Detection P / R / F1            & $1.00\,/\,1.00\,/\,1.00$ \\
Detection TP/FP/TN/FN           & $11/0/3/0$ \\
Patch correctness (5 gates)     & $1.00$ \;(4/4) \\
Stage-G gate-pass rate          & $1.00$ \;(4/4) \\
Abstention soundness            & $1.00$ \;(7/7) \\
Residual sinks (M4)             & $0$ \\
Median latency (M5)             & $\approx 0.0018$\,s \\
\hline
\end{tabular}
\end{table}

\begin{figure}[ht]
\centering
\includegraphics[width=0.78\linewidth]{figures/fig_summary.pdf}
\caption{Controlled-tier guarantees at a glance: perfect detection F1,
perfect patch correctness and perfect abstention soundness, with a
residual-CWE rate of zero. The figure visualises the rows of
Table~\ref{tab:snippet} that bear on the soundness objectives R1--R2.}
\label{fig:summary}
\end{figure}

A per-class breakdown confirms that the perfect aggregate is not hiding a
weak class: CWE-643 is detected and handled $2/2$, CWE-78 $3/3$, CWE-89
$7/7$ and CWE-90 $2/2$, where ``handled'' means either a gate-passing patch
or a sound abstention. One should not over-read fourteen records, since this is a
small, curated corpus by construction, but the tier establishes a clean
existence result: on inputs the engine is designed to handle, it handles
them correctly and certifiably, and on inputs it is not yet designed to
repair, it abstains rather than guesses. That is precisely the behaviour
hypotheses H1 and H2 predict.

\subsection{Project tier: detection}
\label{sec:res-detection}

The realistic tier tells a more textured story, and the detection figures
are the first place the texture shows. Over 849 labelled in-scope OWASP
Benchmark cases, the bundled offline detector scores as in
Table~\ref{tab:bench} and Figure~\ref{fig:detbench}. Aggregate precision is
$0.573$ and recall $0.180$, for an F1 of $0.273$. Recall is, frankly,
modest, and we make no attempt to dress it up: the ruleset is a
conservative, dependency-free default, and the CWE-90 and CWE-643 rows are
identically zero because the bundled taint rules target SQL and command
sinks only. These numbers describe the bundled detector, not the engine, and
under design decision~D5 they are a property of an interchangeable
component. Pairing \irsam{} with a stronger detector would raise the volume
of findings acted upon without touching any soundness guarantee.

\begin{table}[ht]
\centering
\caption{Detection on OWASP BenchmarkJava over 849 labelled in-scope cases,
using the bundled offline ruleset. The figures characterise an
interchangeable detector, not the remediation engine.}
\label{tab:bench}
\begin{tabular}{p{2.4cm}ccccc}
\hline
\textbf{Scope} & \textbf{P} & \textbf{R} & \textbf{F1} & \textbf{TP/FP/TN/FN} & \textbf{n} \\
\hline
Overall  & 0.573 & 0.180 & 0.273 & 79/59/350/361 & 849 \\
CWE-89   & 0.556 & 0.257 & 0.352 & 70/56/176/202 & 504 \\
CWE-78   & 0.750 & 0.071 & 0.130 & 9/3/122/117   & 251 \\
CWE-90   & 0.000 & 0.000 & 0.000 & 0/0/32/27     & 59 \\
CWE-643  & 0.000 & 0.000 & 0.000 & 0/0/20/15     & 35 \\
\hline
\end{tabular}
\end{table}

\begin{figure}[ht]
\centering
\includegraphics[width=0.82\linewidth]{figures/fig_detection_bench.pdf}
\caption{Per-class detection precision, recall and F1 on OWASP
BenchmarkJava with the bundled offline ruleset. Precision is respectable
where the rules apply (CWE-89, CWE-78) but recall is low, and the CWE-90 and
CWE-643 classes are uncovered because the bundled taint rules target SQL and
command sinks only. Detection is an interchangeable component under design
decision~D5.}
\label{fig:detbench}
\end{figure}

\subsection{Project tier: remediation}
\label{sec:res-remediation}

The remediation behaviour at scale is the result that matters, and it is
unambiguous. Across 140 real-world in-scope findings (138 from the OWASP
Benchmark and 2 from java-sec-code, with the Python target dvpwa producing
no in-scope findings under the bundled rules) the pipeline synthesised
\emph{zero} unsound patches and returned a sound abstention in every single
case. Not one false fix slipped through. Table~\ref{tab:abstain} and
Figure~\ref{fig:abstain} break down the Benchmark abstention reasons, and
the breakdown is itself instructive about the structure of real code.

\begin{table}[ht]
\centering
\caption{Abstention reasons over the 138 OWASP Benchmark findings. Every
finding the engine could not repair provably was declined with a precise,
machine-readable reason; no unsound patch was emitted.}
\label{tab:abstain}
\begin{tabular}{p{8.8cm}c}
\hline
\textbf{Reason} & \textbf{Count} \\
\hline
No connection variable in local scope                & 55 \\
Already parameterised / placeholder query (\code{?}) & 54 \\
Non-DML grammar (\code{CallableStatement \{call\}})  & 12 \\
Unsupported backend (\code{java/shell})              & 12 \\
Unsupported quoting syntax                           & 5 \\
\hline
\textbf{Total}                                       & \textbf{138} \\
\hline
\end{tabular}
\end{table}

\begin{figure}[ht]
\centering
\includegraphics[width=0.86\linewidth]{figures/fig_abstention.pdf}
\caption{Distribution of abstention reasons over the 138 OWASP Benchmark
findings. The two dominant reasons, no connection variable in local scope
and a query that is already parameterised, together account for 109 of 138
cases, and both are sound refusals rather than failures: the first reflects
inter-procedural construction the slicer will not follow, the second a sink
that needs no repair.}
\label{fig:abstain}
\end{figure}

Two reasons dominate. Fifty-five findings abstain because no connection
variable is in local scope (the query is assembled across method
boundaries the slicer will not cross) and fifty-four abstain because the
query is already parameterised, carrying a literal \code{?} placeholder and
therefore needing no repair at all. The remaining twenty-nine split across
non-DML grammar, the unsupported shell backend, and quoting syntax outside
the supported fragment. It is worth observing that these are not arbitrary
failures; each corresponds to a precise boundary of the supported fragment
or a property of the code, and each is reported as such. The java-sec-code
findings both abstain with \code{unsupported\_backend:java/shell}, again a
sound refusal. This is the empirical heart of the chapter: at realistic
scale the engine never traded a proof for a guess.

\subsection{Cross-tier consistency and earlier phases}
\label{sec:res-phases}

Earlier evaluation phases, run on larger and differently constructed
corpora, corroborate the two-tier picture. Table~\ref{tab:phases} gathers
the headline figures. A 50-case snippet-level phase achieved $87.5\%$ patch
applicability with zero residual CWE-89; a 21-case multi-language phase
covering Python and JS/TS reached a $100\%$ patch pass rate, again with no
residual sinks; and a 20-case synthetic phase, used for the baseline
comparison below, recorded a functional fix rate (M2) of $0.80$, semantic
equivalence (M3) of $1.00$, a residual-CWE mean (M4) of $0.00$, an
abstention precision (M6) of $1.00$ and a median latency (M5) of three
milliseconds. The figures are not identical across phases, as they should not
be given different corpora, but they are mutually consistent in the
properties that matter: high applicability where the fragment is supported,
zero residual sinks, and sound abstention everywhere else.

\begin{table}[ht]
\centering
\caption{Headline figures from the principal evaluation phases. ``Applic.''
is patch applicability (M1); ``Resid.'' is residual-CWE count or mean (M4).}
\label{tab:phases}
\begin{tabular}{p{4.4cm}cccc}
\hline
\textbf{Phase / corpus} & \textbf{Cases} & \textbf{Applic.} & \textbf{Resid.} & \textbf{Note} \\
\hline
Snippet-level (P2)       & 50 & $0.875$ & $0$ & 35/40 patched, 5/10 abstain \\
Multi-language (P4)      & 21 & $1.000$ & $0$ & Python + JS/TS \\
Synthetic (P5)           & 20 & $0.800$ & $0$ & M2$=0.80$, M6$=1.00$ \\
Controlled snippet       & 14 & $1.000$ & $0$ & 4/4 patched, 7/7 abstain \\
Project (Benchmark)      & 138 & N/A & $0$ & 138/138 sound abstentions \\
\hline
\end{tabular}
\end{table}

\subsection{Comparison with baselines}
\label{sec:res-baselines}

Against the six baselines on the 20-case synthetic evaluation, \irsam{}
separates cleanly on the metric that counts. Figure~\ref{fig:m2} plots the
functional fix rate. \irsam{} fixes $0.80$ of cases; the strongest baseline,
Semgrep's autofix, manages $0.20$; the CodeQL-plus-GPT-4 pipeline reaches
$0.15$; and the three learning-based repairers, VulRepair, SeqTrans and
ChatRepair, each land at $0.05$, with the zero-shot LLM at $0.00$. The
$4\times$ margin over the nearest competitor is large, but the more telling
comparison is multi-dimensional. Figure~\ref{fig:radar} overlays five
metrics at once. \irsam{} holds applicability, fix rate, semantic
equivalence and abstention precision high while keeping residual sinks at
zero; the autofix and LLM baselines either fix little or, where they do
emit edits, do so without the abstention discipline that keeps M6 high.

\begin{figure}[ht]
\centering
\includegraphics[width=0.74\linewidth]{figures/fig_m2_baselines.pdf}
\caption{Functional fix rate (M2) on the 20-case synthetic evaluation.
\irsam{} repairs $0.80$ of cases against $\le 0.20$ for every baseline,
while, as the other metrics confirm, holding residual sinks at zero and
abstention precision at one. The margin reflects the value of provable
repair on the supported fragment, not a broader coverage claim.}
\label{fig:m2}
\end{figure}

\begin{figure}[ht]
\centering
\includegraphics[width=0.66\linewidth]{figures/fig_metrics_radar.pdf}
\caption{Multi-metric comparison of \irsam{} against three representative
baselines (M1 applicability, M2 fix rate, M3 equivalence, the complement
$1-$M4 of residual density, and M6 abstention precision; larger is better on
every axis). \irsam{} dominates on the fix-rate and equivalence axes while
matching the best baseline on residual density and abstention precision.}
\label{fig:radar}
\end{figure}

A fair reading of these plots requires a caveat we are happy to supply. The
synthetic corpus is weighted toward the fragment \irsam{} supports, so the
comparison measures how well each tool does \emph{when a sound repair
exists}, not how broadly each tool applies. Generative baselines would
out-cover \irsam{} on, say, memory-safety bugs the present engine does not
address at all. The claim here is narrower and, we think, more defensible:
within injection remediation, prove-or-abstain yields both a higher fix rate
and a zero false-fix rate, which is the combination security practice
actually wants.

\subsection{Iteration and defects fixed}
\label{sec:iteration}

Honesty about the evaluation requires reporting that it was iterative, and
that running the harnesses surfaced one detection gap and two genuine gate
defects before the figures above stabilised. We recount them because the
episode illustrates the value of the multi-gate design. First, purely
syntactic detection rules returned zero findings on the OWASP Benchmark,
which assembles its queries across statements; adding two taint-mode rules
raised Benchmark findings from 0 to 138 and enabled Table~\ref{tab:bench},
while the 35 detector, gate and ingest unit tests continued to pass. Second,
the compile gate ran \code{javac} unconditionally and therefore failed every
Python patch; it now selects the language-appropriate compiler. Third, the
differential gate flagged a patched-statement error as a regression even
when the original statement also errored on the fixture (for instance a
query against a table absent from the fixture schema), and it now records a
regression only against a working baseline. Before these fixes, the four
repairable SQL snippets synthesised correct patches but failed stage~G, a
gate-pass rate of $0/4$; afterward the rate is $4/4$ with zero residual
sinks, and the full 336-test suite remained green throughout. The defects
were in the \emph{gates}, not in the soundness logic, which is reassuring:
the engine never accepted an unsound patch even while the gates were
miscounting.

\subsection{Failure-mode analysis}
\label{sec:res-failures}

A failure-mode study over the synthetic corpus rounds out the picture and,
usefully, locates the failures that do occur outside \irsam{}.
Table~\ref{tab:failures} lists the three dominant modes. The largest,
accounting for roughly a third of observations, is an uncategorised bucket
dominated by Semgrep's autofix; the next two, together another half, are
model-generation failures and hallucinated APIs, and both are attributed to
the generative baselines rather than to \irsam{}. In other words, the
recurrent failure modes in the comparison are properties of guess-based
repair: an LLM inventing an API that does not exist is the canonical
example, not a property of the prove-or-abstain engine. We read this as indirect
but pointed support for the central design premise.

\begin{table}[ht]
\centering
\caption{Top failure modes observed across the baseline comparison. The two
generative modes (model-generation failure, hallucinated API) are properties
of guess-based baselines, not of \irsam{}.}
\label{tab:failures}
\begin{tabular}{p{5.4cm}ccp{3.0cm}}
\hline
\textbf{Failure mode} & \textbf{Count} & \textbf{\% obs.} & \textbf{Dominant tool} \\
\hline
Uncategorised            & 46 & 32.9 & Semgrep-autofix \\
Model-generation failure & 40 & 28.6 & VulRepair \\
Hallucinated API         & 28 & 20.0 & LLM zero-shot \\
\hline
\end{tabular}
\end{table}

\subsection{Revisiting the hypotheses}
\label{sec:res-hypotheses}

We can now read the evidence back against Section~\ref{sec:hypotheses}.
Hypothesis H1, that provable repair is feasible on a useful fragment, is
supported: on the controlled tier every repairable case was fixed with a
re-checkable certificate, and the synthetic tier sustained an M2 of $0.80$.
Hypothesis H2, that abstention is sound and common rather than rare, is
strongly supported: across 140 real findings the engine abstained soundly
$140/140$, with abstention precision $1.00$ wherever measured, and the
abstention reasons in Table~\ref{tab:abstain} are concrete and
code-grounded. Hypothesis H3, that soundness need not collapse the fix rate,
is supported within scope: the $4\times$ margin over the nearest baseline in
Figure~\ref{fig:m2} shows that insisting on provable repair, on the
supported fragment, leaves the functional fix rate substantially \emph{above}
that of tools willing to guess. None of this contradicts the honest
limitation that the supported fragment is narrower than the space of all
injection sinks, a point we take up next.


% =====================================================================
\section{Discussion}
\label{rem:discussion}
% =====================================================================

Having reported what the engine does, we step back to interpret it: what the
results mean, where the approach is weak, which threats might undermine the
conclusions, how the work sits beside the related literature, and what we
would build next. We try to be our own harshest critic here, since a
soundness claim invites, and deserves, scrutiny.

\subsection{Interpretation}
\label{sec:disc-interpret}

The central finding admits a one-sentence summary: on the fragment it
supports, \irsam{} repairs injection vulnerabilities with a machine-checkable
certificate, and everywhere else it abstains with a precise reason, never
emitting an unsound patch across 140 real-world findings. The significance is
less the headline fix rate than the \emph{shape} of the guarantee. A
practitioner integrating \irsam{} into a pipeline obtains a tool whose
accepted patches can be trusted without re-review of the security property,
and whose silences are informative rather than opaque. That combination, a
positive output that is certified and a negative output that is
explained, is exactly what a guess-based repairer cannot offer, because it
has no notion of declining when it cannot prove. The $4\times$ fix-rate
margin in Figure~\ref{fig:m2} should therefore be read together with the
zero false-fix rate: the method does not buy soundness by becoming timid, at
least not within its scope.

A second interpretive point concerns the project-tier abstentions. That
$109$ of $138$ Benchmark findings abstain for just two reasons: no
local connection variable, or an already-parameterised query. This is not a
disappointment but a measurement. It tells us that real injection-prone code
is dominated by inter-procedural query construction and by sinks that are
already safe, which in turn tells us precisely where engineering effort would
move the coverage needle. An evaluation that merely reported ``low coverage''
would hide this; the machine-readable abstention reasons expose it.

\subsection{Limitations}
\label{sec:disc-limits}

We are candid about the boundaries. The supported repair fragment is
$\mathrm{SQL}_0$ for the SQL interpreter, with full A--G backends only for
Java and Python; the shell, LDAP and XPath interpreters are modelled and
abstained rather than repaired, despite shipping binders. Query construction
that crosses more than one method boundary defeats the slicer by design.
Detection, as Table~\ref{tab:bench} makes plain, is limited by the
conservative bundled ruleset, so the volume of findings acted upon is lower
than a production detector would supply. And the Lean development, while free
of \code{sorry} in its proved core, still carries eight per-framework
wire-protocol lemmas discharged only by placeholder, so the end-to-end
mechanised guarantee does not yet extend to every catalog. None of these is
fatal to the thesis: each is a boundary of \emph{scope}, not a soundness
hole, but together they bound the generality of the claims, and we would
rather state them than have a reader discover them.

\subsection{Threats to validity}
\label{sec:disc-threats}

Following standard practice we separate three kinds of threat. The
\emph{construct} threat is whether our metrics measure what we claim: M2
counts a fix only when the sink is removed and benign behaviour preserved,
and the structural gate certifies the security property directly, so we
believe the construct is tight, though M2 inherits whatever blind spots the
differential fixture has. The \emph{internal} threat is whether the results
are artefacts of the harness; the iterative defect-fixing of
Section~\ref{sec:iteration} is the clearest internal risk, and we mitigate it
by reporting the defects openly and by noting that the 336-test suite
remained green across the changes. The \emph{external} threat is
generalisation: the controlled and synthetic corpora are small and curated,
and the single project at scale is the OWASP Benchmark, which is synthetic by
nature. The two genuine open-source targets (java-sec-code, dvpwa) are small.
We therefore make no claim about behaviour on large proprietary codebases,
and we flag corpus breadth as the first thing future work should expand.

\subsection{Critical view of related work}
\label{sec:disc-related}

Set against the related work surveyed in Chapter~\ref{chap:state-of-the-art}, \irsam{}
occupies a deliberately narrow and deliberately rigorous niche. Learning-based
repairers such as VulRepair \cite{fu2022vulrepair}, SeqTrans
\cite{chi2023seqtrans} and the conversational approach of ChatRepair
\cite{xia2023chatrepair} cover a far wider range of bug classes than we do,
and that breadth is a genuine virtue; but, as Figure~\ref{fig:m2} and the
hallucinated-API failure mode in Table~\ref{tab:failures} illustrate, they
offer no guarantee that a proposed edit is correct or even compiles, and they
provide no principled way to decline. Pattern-based autofix, exemplified by
Semgrep \cite{semgrep2024} and the CodeQL-assisted pipelines now appearing in
practice \cite{githubautofix2024}, is fast and deterministic but inherits the
imprecision of its patterns and, again, lacks an abstention discipline. Our
position is not that these tools are wrong (they answer a broader question
than we do), but that for the specific, high-stakes class of injection
remediation, a prove-or-abstain engine supplies a guarantee the others
structurally cannot. One may reasonably argue the opposite trade-off for
contexts where coverage matters more than certified correctness; we would
simply observe that security remediation is rarely such a context.

\subsection{Future work}
\label{sec:disc-future}

Several extensions follow directly from the limitations. The most immediate
is to register host rewriters for the shell, LDAP and XPath interpreters,
whose binders and SIG schemas already exist, turning sound abstentions
into sound repairs with no change to the trusted core. The second is to widen
the grammar fragment beyond $\mathrm{SQL}_0$ toward stored-procedure calls and
the quoting constructs that Table~\ref{tab:abstain} shows are common. The
third is to pair the engine with a production-grade detector, which the
detector-agnostic ingest layer already accommodates, so that the recall
ceiling of Table~\ref{tab:bench} is lifted without touching any guarantee.
The fourth, and most foundational, is to discharge the eight remaining Lean
lemmas by supplying an equational model of the JDBC and DB-API wire protocols,
extending the mechanised guarantee end to end across every catalog. Finally,
following the slicer's dominant abstention reason, an inter-procedural slice
that crosses bounded method boundaries would directly address the largest
single class of real-world abstentions. We regard the first and fourth of
these as the highest-value next steps.


% =====================================================================
\section{Conclusion}
\label{rem:conclusion}
% =====================================================================

This chapter has developed \irsam{}, an engine for automated remediation of
injection vulnerabilities whose defining commitment is to repair only when it
can prove the repair sound, and to abstain, providing a precise,
machine-readable reason, whenever it cannot. We began from the observation
that automated remediation has historically faced an uncomfortable choice
between guessing, with the attendant risk of false fixes, and doing nothing;
the prove-or-abstain stance dissolves that dilemma by making ``do nothing''
an informative, sound action rather than a failure.

The contribution is concrete on several fronts, and it is worth tying each
back to the objectives of Section~\ref{sec:objectives}. We formalised the
remediation problem and the soundness criterion (R1), realising it as the
seven-stage A--G pipeline of Figure~\ref{fig:pipeline} that lifts a detector
finding through slicing, reconstruction, intent parsing, binding, synthesis
and a five-gate validator. We mechanised the central soundness argument in
Lean~4 (R2), closing twenty theorems with no \code{sorry} in the proved core,
including the principal parameterisation and audit-soundness results. We made
the engine detector-agnostic (R4) through a unified finding schema, and
externalised its safe-API knowledge into eleven schema-validated binder
catalogs (R3) so that coverage extension is a data change rather than a code
change. And we evaluated the system end to end against real tools (R5--R6):
perfect detection, patch correctness and abstention soundness on the
controlled tier; zero unsound patches across 140 real-world findings, every
one a sound abstention; and a functional fix rate roughly four times the
nearest of six baselines on the supported fragment.

We have been equally clear about what the work does not yet do. The repair
fragment is bounded, three interpreters are modelled but not yet rewritten,
the bundled detector's recall is modest, and eight wire-protocol lemmas
remain to be discharged in Lean. These are boundaries of scope, and each
maps to a concrete next step rather than to a flaw in the guarantee. The
overarching lesson, we believe, generalises beyond injection: when an
automated tool's output carries security weight, the discipline of proving
before patching, and of abstaining audibly when no proof is in reach, is
not a limitation to be apologised for but the very property that makes the
tool's accepted output worth trusting. It is on that conviction that the
remediation contribution of this thesis rests, and from which its future
extensions naturally proceed.